\documentclass[10pt,a4paper]{article}

\usepackage[a4paper,left=0.30in,right=0.30in,top=0.27in,bottom=0.27in]{geometry}
\usepackage{fontspec}
\setmainfont{Arial}
\usepackage[english]{babel}
\usepackage{enumitem}
\usepackage{xcolor}
\usepackage{hyperref}
\usepackage{titlesec}
\usepackage{microtype}

\definecolor{primary}{HTML}{2463A5}
\definecolor{textdark}{HTML}{303030}
\definecolor{textlight}{HTML}{666666}
\definecolor{rulecolor}{HTML}{C7C7C7}
\hypersetup{colorlinks=true,urlcolor=primary,linkcolor=primary}

\pagestyle{empty}
\setlength{\parindent}{0pt}
\setlength{\parskip}{0pt}
\linespread{0.78}
\urlstyle{same}

\titleformat{\section}
  {\normalsize\bfseries\color{primary}}
  {}{0pt}{\MakeUppercase}
  [\vspace{-0.3em}\color{rulecolor}\rule{\linewidth}{0.45pt}]
\titlespacing*{\section}{0pt}{0.34em}{0.16em}

\setlist[itemize]{
  label=\textbullet,
  leftmargin=1.1em,
  itemsep=0.04em,
  topsep=0.07em,
  parsep=0pt,
  partopsep=0pt
}

\newcommand{\entry}[3]{%
  \vspace{0.08em}%
  \noindent\textbf{\color{textdark}#1}\hfill{\color{textlight}#3}\par
  \noindent{\small\color{primary}#2}\par
}

\newcommand{\project}[3]{%
  \vspace{0.06em}%
  \noindent\textbf{\color{textdark}#1}\hfill{\color{textlight}#3}\par
  \noindent\begingroup\small\color{textdark}#2\par\endgroup
}

\begin{document}

% Full-width header and summary keep the reading order clear for ATS parsers.
{\fontsize{22}{24}\selectfont\textbf{\color{textdark}Galaba Vamsi}}\par
\vspace{0.08em}
{\small\color{primary}Software Engineer | AI/ML | Computer Vision | Systems}\par
\vspace{0.14em}
{\small\color{textlight}
\href{tel:+917483047846}{+91 7483047846} $\vert$
\href{mailto:galabavamsi12@gmail.com}{galabavamsi12@gmail.com} $\vert$
\href{https://galabavamsi.github.io/portfolio}{galabavamsi.github.io/portfolio}\\
\href{https://www.linkedin.com/in/galabavamsi/}{linkedin.com/in/galabavamsi} $\vert$
\href{https://github.com/Galabavamsi}{github.com/Galabavamsi}}

\section*{Professional Summary}
{\small\color{textdark}
Mechatronics undergraduate at IIT Bhilai focused on software engineering, AI/ML, computer vision, robotics, and GPU systems. Builds full-stack products, open-source tooling, simulators, and visual applications with Python, TypeScript, React, PyTorch, CUDA, and cloud platforms.
}

\vspace{0.12em}
\noindent\begin{minipage}[t]{0.645\textwidth}
\vspace{0pt}
\small

\section*{Experience}

\entry{Open-Source Contributor - HFlow}{Hebbian Robotics (YC S26) | \href{https://github.com/Hebbian-Robotics/hflow}{HFlow repository}}{2026}
\begin{itemize}
  \item Delivered 7 maintainer-reviewed merged PRs addressing 7 tracked issues in HFlow's robotics and Physical AI SDK.
  \item \href{https://github.com/Hebbian-Robotics/hflow/pull/598}{PR \#598} / \href{https://github.com/Hebbian-Robotics/hflow/issues/597}{issue \#597}: refused sources with multiple channels on one topic before canonical episode creation, added doctor diagnostics, and bumped \texttt{TRANSFORM\_BEHAVIOR\_VERSION} to 10.
  \item \href{https://github.com/Hebbian-Robotics/hflow/pull/350}{PR \#350} / \href{https://github.com/Hebbian-Robotics/hflow/issues/296}{issue \#296}: implemented paginated LeRobot tree traversal using \texttt{Link: rel="next"} headers, preserved request headers, deduplicated paths, rejected unsafe cross-origin URLs, and prevented pagination loops.
  \item \href{https://github.com/Hebbian-Robotics/hflow/pull/353}{PR \#353} / \href{https://github.com/Hebbian-Robotics/hflow/issues/292}{issue \#292}: resolved LeRobot video-cache collisions by including camera, chunk, and file indexes in cache identity and bumping the converter version.
  \item \href{https://github.com/Hebbian-Robotics/hflow/pull/360}{PR \#360} / \href{https://github.com/Hebbian-Robotics/hflow/issues/314}{issue \#314}: centralized safe ingest URI parsing across the CLI, server, and SDK and added validation for blank, absolute, and parent-escaping paths.
  \item \href{https://github.com/Hebbian-Robotics/hflow/pull/367}{PR \#367} / \href{https://github.com/Hebbian-Robotics/hflow/issues/365}{issue \#365}: benchmarked 3 cold 1080p30 checks (900 frames/check; 4.897 s median) and separated transform timing/provenance; \href{https://github.com/Hebbian-Robotics/hflow/pull/400}{PR \#400} / \href{https://github.com/Hebbian-Robotics/hflow/issues/399}{issue \#399}: refused LeRobot depth videos before conversion to preserve 12-bit depth semantics.
\end{itemize}

\entry{Open-Source Contributor - Inspect Robots}{Robocurve (YC Summer 2026) | \href{https://github.com/robocurve/inspect-robots}{Inspect Robots repository}}{Sep 2026 - Present}
\begin{itemize}
  \item Contributed 3 PRs to Inspect Robots: \href{https://github.com/robocurve/inspect-robots/pull/476}{\#476} merged configurable retry policy across agent integrations; \href{https://github.com/robocurve/inspect-robots/pull/475}{\#475} remains open for preserving unknown git provenance; and \href{https://github.com/robocurve/inspect-robots/pull/434}{\#434} remains open for a first-party W\&B logging sink plugin.
\end{itemize}

\entry{Research Intern - 6G RIS Simulator}{IIT Bhilai | Advisor: Dr. Arzad Alam}{May 2025 - Aug 2025}
\begin{itemize}
  \item Built a zero-budget, open-source Python simulator with USRP-style interfaces for line-of-sight and non-line-of-sight wireless environments; RIM-SIM v2 includes 6 paper-grounded scenarios, a 20 Hz WebSocket feed, and 103 passing tests. \href{https://github.com/Galabavamsi/RIM-SIM-V2}{[GitHub]}
  \item Designed and validated a 2-layer reconfigurable intelligent surface (RIS) PCB workflow for self-powered RF energy harvesting; contributed to 3 related publications (1 as first author and 2 as co-author).
\end{itemize}

\entry{Marketing Intern - Campus CEO Project}{Swiggy Ltd}{May 2024 - Jul 2024}
\begin{itemize}
  \item Managed a campaign with 15 social media posts; contributed to campus marketing initiatives and execution.
\end{itemize}

\section*{Projects}

\project{Human Slop - Co-Founder}{Built and deployed an anti-AI web and Android platform with React, PostgreSQL, and AWS; implemented hardware-bound biometric authentication and real-time typing forensics. \href{https://humanslop.in}{[Website]} | \href{https://play.google.com/store/apps/details?id=com.humanslop.app}{[App]} | \href{https://github.com/Galabavamsi/human-slop}{[GitHub]}}{2026 - Present}

\project{Charter - Agentic Commerce}{Built a TypeScript/Fastify/Postgres/Vite + React agentic-commerce platform with 3 interfaces (chat, Vapi voice, MCP/HTTP) and 10 MCP tools over a shared commerce core; implemented grounded catalog and policy tools, frozen INR quote state, tenant-scoped orders, and RLS. \href{https://github.com/Galabavamsi/charter}{[GitHub]}}{}

\project{Pneumatic Sorting Digital Twin}{Built an open-source Unity digital twin for modular electropneumatic automation with offline sorting and stamping kits, HMI flows, telemetry, and replay. \href{https://github.com/Galabavamsi/pneumatic-sorting-digital-twin}{[GitHub]}}{}

\project{Arista Networks Wi-Fi Optimization System - Inter IIT 5th Rank}{Built a client-aware Wi-Fi RRM system in NS-3 using 3 RF metrics; classified 100,000 I/Q samples from 5 ms captures into 16 classes with STFT, DBSCAN, and Dual-CNN inference; simulated 5 access points and 50 client stations using IEEE 802.11k/v. \href{https://canva.link/puoqbe4okoid8rb}{[Slides]}}{2025}

\project{Electron-GNN - ML-Accelerated Quantum Spectroscopy}{Built a V3 two-tower GATv2 model to predict molecular absorption spectra from geometry; evaluated on 2 molecules with 39 and 55 target peaks. \href{https://github.com/Galabavamsi/Electron-GNN}{[GitHub]}. Faculty advisor: \href{https://www.iitbhilai.ac.in/index.php?pid=soumajit}{Dr. Soumajit Pramanik}.}{2026}

\project{Open-PyFX - GPU-Accelerated Video Effects}{Built a Python and GLSL video effects suite with GPU rendering for CRT emulation, ASCII edge detection, retro graphics, and 3D LUT pipelines; supports 3840x2160 export and up to 16x zoom. \href{https://github.com/Galabavamsi/open-pyfx}{[GitHub]}}{2025}

\project{Lunar Crater Detection - DEM Terrain Analysis}{Built a computer vision pipeline for crater extraction from Digital Elevation Models; achieved an F1 score of 0.839 and approximately 4x preprocessing speedup using CuPy/CUDA; the repository includes 77 tests. \href{https://github.com/Galabavamsi/Moon-Crater-Detection-using-DEM}{[GitHub]}}{2025}

\project{AntennaNet - Inverse Electromagnetic Design}{Developed KD-tree nearest-neighbour spectral anchoring for inverse electromagnetic design across 144D search spaces using Python, MEEP, and OpenMPI; indexed approximately 2,000 spectra and a 1,200-sample, 12x12-pixel dataset. \href{https://github.com/Galabavamsi/Antenna-Net}{[GitHub]}}{2025}

\project{Amazon ML Challenge - Top 3.5\%}{Multimodal price prediction with Qwen 2-VL 7B and a ViT/RoBERTa two-tower model; 75,000 data pairs; SMAPE 49.0.}{2025}

\end{minipage}\hfill
\begin{minipage}[t]{0.325\textwidth}
\vspace{0pt}
\footnotesize
\raggedright

\section*{Skills}
{\color{textdark}
\textbf{Languages:} Python, C/C++, TypeScript, JavaScript, SQL, GLSL, CUDA.\par
\textbf{AI/ML:} PyTorch, Graph Neural Networks (GATv2), Safe Reinforcement Learning, Transformers, OpenCV, CuPy, NumPy.\par
\textbf{Software and Systems:} React, React Native, Flutter, FastAPI, Node.js, REST APIs, WebSockets, PostgreSQL, AWS, Supabase, Docker, ROS 2, Gazebo, MEEP, OpenMPI, SDR/USRP.\par
\textbf{Developer Tools:} Git, Linux, GitHub Actions, CI/CD, pytest, Ruff, ty, static analysis, regression testing.
}

\section*{Education}
\entry{B.Tech. in Mechatronics}{Indian Institute of Technology Bhilai}{Aug 2023 - May 2027}
{\color{textdark}\textbf{Coursework:} Data Structures and Algorithms, Object-Oriented Programming, Deep Learning for Computer Vision, AI and Machine Learning, Digital Image Processing, Optimal Control.}

\section*{Publications}
{\color{textdark}
\textbf{First author:} \textit{An Open Emulator for Smart Radio Environments}. \href{https://drive.google.com/file/d/12oDPsflaUXKnSjfA3247HgFDDf4bJkCP/view?usp=drive_link}{[Link]}\par
\textbf{Co-author:} \textit{Experience with RF Energy Harvesting-Driven Self-Powered RIS}, IEEE INDICON 2025. \href{https://ieeexplore.ieee.org/abstract/document/11392908/}{[Link]}\par
\textbf{Co-author:} \textit{A systematic literature review on simulation models and deployments for reconfigurable intelligent surfaces}. \href{https://www.sciencedirect.com/science/article/pii/S1570870526001964}{[Link]}
}

\section*{Achievements}
{\color{textdark}
\textbf{1st Place:} IIT Bhilai Web Development Competition.\par
\textbf{5th Rank:} Arista Networks Wi-Fi Challenge, Inter IIT Tech Meet 14.0.\par
\textbf{Top 3.5\%:} Amazon ML Challenge.\par
\textbf{YC Startup School:} Human Slop accepted.\par
\textbf{Top 10 Finalist:} Toyota Hackathon.
}

\end{minipage}

\end{document}
